\begin{topic}[id=ics_energy_attacks,
             difficulty={Mittel},
             type={Sicherheitsanalyse + Fallbezug},
             prerequisites={Grundlagen IT-Sicherheit; Interesse an industriellen Systemen und Kritischen Infrastrukturen},
             practice={optional},
             owner={Dozententeam},
             updated={2026-04-09},
             tags={ ICS, SCADA, KRITIS, IT/OT, Zoning, Fernwartung, Monitoring }]{Cyberangriffe auf Industrie \& Energie}

\TopicDescription{In Industrie/Energie wirkt Security direkt in die physische Welt. Typische Pfade: unsichere Fernwartung, ungepatchte Komponenten, schwache Segmentierung. Das Thema verknüpft reale Vorfälle mit Gegenmaßnahmen: Netzwerk-Zoning, Whitelisting, Härtung von Protokollen/Services, OT-sichtbares Monitoring und geübtes Incident-Management. Fokus ist die IT/OT-Schnittstelle: Zusammenarbeit, Telemetrie und sichere Änderungen im Betrieb.}

\TopicLeitfragen{\item Welche Zonen/Conduits sind in der Praxis sinnvoll?
\item Wie wird Fernwartung sicher geplant und überwacht?
\item Wie verbindet man IT- und OT-Monitoring wirksam?
}
\TopicScope{\item Keine Tool-/Vendor-Vergleiche im Detail.
\item Keine regulatorische Tiefenprüfung.
}
\TopicPractice{Skizze eines Zoning-Konzepts inkl. Fernwartung und Patch-Prozess.}

\TopicKeywords{ICS, SCADA, KRITIS, IT/OT, Zoning, Fernwartung, Monitoring}

\nocite{cisaICS,enisaICS}
\end{topic}
